% ============================================================================
% VI. DISCUSSION (REVISED - FINAL)
% ============================================================================
\section{Discussion}
\label{sec:discussion}

% ----------------------------------------------------------------------------
\subsection{Key reasons for the CAE-induced improvements}
\label{sec:discussion_why}

This wavelength selection approach showed that using fewer bands yields better results than using the full 192 bands(Table~\ref{tab:main_results}). Accuracy did not decrease monotonically with band reduction; instead, an optimal operating point existed where sufficient discriminative information was retained while noise was maximally suppressed.
Beyond this optimum, further reduction traded accuracy for efficiency, a favorable exchange for resource-constrained applications.

\subsubsection{Noise Removal Through Intelligent Selection}
Hyperspectral data inherently contains redundancy and noise.
Spectral bands with low signal-to-noise ratios, bands capturing uninformative variance (e.g., illumination artifacts), and bands with high inter-class overlap all contribute to the full spectrum.
A classifier operating on all 192 bands had to navigate this noise, which could degrade decision boundaries.

By selecting only wavelengths with high discriminative value, the method effectively performed dimensionality reduction that removed uninformative or harmful information.
The perturbation analysis identified wavelengths that the autoencoder considered important for faithful reconstruction, wavelengths that captured meaningful spectral structure rather than noise.

\subsubsection{Joint 4D Representation}
Traditional band selection methods process each spectral band independently, assuming that informativeness can be assessed from individual band statistics.
For ME-HSI data, where discriminative patterns emerge from excitation-emission \emph{combinations}, this assumption failed.
The parallel-branch architecture with feature averaging enabled the network to learn cross-excitation correlations---patterns that manifested only when comparing fluorescence responses across excitation wavelengths.
The latent space encoded these joint patterns, capturing information that per-band analysis would miss.

\subsubsection{Perturbation as Learned Attribution}
The perturbation analysis operated on what the autoencoder had learned, not on raw data statistics.
By measuring reconstruction sensitivity to latent perturbations, wavelengths that the network considered essential for accurate reconstruction were identified.
This is fundamentally different from variance-based band selection, which might prioritize noisy bands with high variance but low discriminative value.
The autoencoder's compression objective acted as an implicit importance filter: wavelengths that did not contribute to reconstruction fidelity were effectively ignored.

% ----------------------------------------------------------------------------
\subsection{Cross-Dataset Generalization}
\label{sec:discussion_cross_dataset}

The framework was evaluated on two ME-HSI fluorescence datasets with deliberately different characteristics: Lichens (4 morphological classes, 192 bands) and Collagen Sponges (3 crosslinker-concentration classes, 158 bands).
On both, the method achieved above-baseline performance with substantial data reduction: 95.2\% at $K=80$ on Lichens ($+7.0$ pp) and 85.6\% at $K=30$ on Collagen Sponges ($+5.8$ pp).
The two datasets cover different sample chemistries (lichen autofluorescence vs.\ collagen autofluorescence), different acquisition grids (8 excitations $\times$ variable emission for Lichens, 6 excitations for Collagen Sponges), and different class structures (morphological vs.\ chemical concentration), yet a single hyperparameter family produces top-tier selections on each.
The dimension-scoring strategy proves to be the only meaningful per-dataset parameter: Lichens prefers PCA-$k{=}3$ while Collagen Sponges prefers variance-$k{=}1$; the rest of the pipeline runs identically.

% ----------------------------------------------------------------------------
\subsection{Multi-Classifier Robustness}
\label{sec:discussion_multi_classifier}

The Collagen Sponges multi-classifier evaluation (Section~\ref{sec:results_pepsin}) confirms that the selected bands carry information that is portable across classifier families.
LDA, KNN, SVM (linear and RBF), MLP, Random Forest, and Gradient Boosting Machines were all evaluated on the same selection; nine of ten classifiers benefited.
The strongest individual result was LDA at 92.5\% accuracy on $K=50$ bands (vs.\ 92.8\% on all 158), demonstrating that the selection retains nearly all of the linearly separable structure of the original spectrum.
This counters the concern that selection results may be classifier-specific: the chosen wavelengths are intrinsically informative, not merely tuned to KNN's distance metric.

% ----------------------------------------------------------------------------
\subsection{Robustness Analysis Insights}
\label{sec:discussion_robustness}
The search space of $\binom{192}{13} \approx 3.6 \times 10^{18}$ possible 13-band combinations is far too large for exhaustive search.
That the learned selection outperformed all 10,000 random samples---with a $12.3\sigma$ separation from the random mean ($p \ll 10^{-20}$)---demonstrated that the perturbation-based attribution genuinely identified discriminative wavelengths rather than benefiting from chance.

% ----------------------------------------------------------------------------
\subsection{Object-Level Performance Patterns}
\label{sec:discussion_objects}
The object-level analysis (Table~\ref{tab:objects}) revealed that difficult-to-classify specimens benefited most from wavelength selection.
This suggests that the full spectrum contains wavelengths \emph{detrimental} to certain specimens---likely bands where intra-class variance is high or spectral signatures overlap with other classes.
Removing these bands preferentially improves challenging cases while maintaining performance on straightforward ones.
This property can be valuable since the challenging specimens are often the most diagnostically important (e.g., early-stage lesions with subtle spectral signatures).

The accompanying 54\% reduction in performance variance across specimens (Table~\ref{tab:objects}) makes classification more predictable and less prone to bias, which can be important for quality control applications.

% ----------------------------------------------------------------------------
\subsection{PCA vs.\ Variance Dimension Selection}
\label{sec:discussion_pca}

The substantial gap between PCA and variance-based selection reflects a difference in how these methods analyze the latent space.
PCA identifies orthogonal directions in the autoencoder's learned representation, capturing structured variation through a linear combination of dimensions.
Variance-based selection ranks individual dimensions independently, missing these interactions, hence its lower performance ceiling and need for more bands to reach it.

% ----------------------------------------------------------------------------
\subsection{The Role of Diversity}
\label{sec:discussion_diversity}

The selection with MMR parameter of $\lambda = 1$ underperformed the balanced setting ($\lambda = 0.5$) used throughout the experiments.
This result reflected the redundancy structure of fluorescence spectra.

High-influence wavelength pairs often clustered together: if one emission wavelength at 480\,nm was highly influential, adjacent wavelengths at 470\,nm and 490\,nm likely showed similar influence due to spectral overlap.
Selecting all three provided little additional information but occupied three of the $K$ available selection positions.
MMR's diversity term penalized spectrally similar selections, encouraging the algorithm to ``spread out'' across the EEM space.

The $\lambda = 0.5$ balanced these considerations: wavelengths had to be both relevant (high influence) and complementary (spectrally distinct), mirroring findings in information retrieval where query-relevant layers are more useful when they also provide novel information~\cite{carbonell1998use}.

% ----------------------------------------------------------------------------
\subsection{Limitations}
\label{sec:discussion_limitations}

% \subsubsection{Dataset Scope}
% Evaluation focused on a single primary dataset (lichen fluorescence) with four classes.
% While results were encouraging, validation across diverse specimen types, imaging conditions, and larger class vocabularies would strengthen generalizability claims.
% Future work should evaluate the framework on datasets from different application domains, such as biomedical tissue imaging or agricultural produce inspection.

\subsubsection{Computational Cost and Amortization}
The complete pipeline---autoencoder training plus perturbation analysis across all configurations---took approximately 2 hours on the hardware used for the Lichens dataset (30 minutes for autoencoder training, 90 minutes for perturbation analysis sweep).
Approximately 75\% of total compute is spent in the perturbation step, with the remainder split between data loading, autoencoder optimization, and MMR selection.
This overhead must be weighed against the downstream amortization: once the $K$-band selection is identified, every subsequent classification, segmentation, or anomaly-detection task on the same sample type runs on $K$ features rather than 192, yielding a $192/K \approx 21\times$ speedup at $K=9$ with proportional reductions in storage, transmission, and sensor-array complexity.
For applications with non-negligible downstream compute (per-pixel inference at deployment, large-scale archival), this trade is favorable; for one-shot exploratory analyses on a single sample, the overhead may be prohibitive.

In practice, the perturbation phase is also embarrassingly parallel across (latent-dimension, magnitude) pairs, and an optimized batched implementation reduces wall-clock time by roughly 3$\times$ at the cost of more peak memory.
A faster variant is left to future work but does not affect any reported numerical result.

\subsubsection{Black-Box Attribution}
Perturbation analysis revealed \emph{which} wavelengths the network considered important but not \emph{why}.
Physical interpretation (e.g., attributing high-influence wavelengths to specific fluorophores such as chlorophyll or lignin) requires additional domain expertise \cite{meysurova2021optical} and cannot be derived from the method alone.
Incorporating such knowledge bases could enable more interpretable wavelength attributions.

\subsubsection{Threshold Selection}
The number of bands to select ($K$) is user-specified.
While the results showed that 9 bands exceeded baseline performance for this dataset, choosing an appropriate $K$ for a different application requires domain knowledge or cross-validation.
Automated threshold selection based on performance curves remains an open problem.

% ----------------------------------------------------------------------------
\subsection{Practical Implications}
\label{sec:discussion_practical}

\subsubsection{Sensor Design Guidelines}
The one-per-excitation pattern observed in the 9-band configuration (Section~\ref{sec:results_wavelengths}) directly informs sensor design: coverage across the full excitation range is essential, but only a single emission band per excitation can be needed for effective discrimination.
The 95\% data reduction translates to proportional reductions in sensor complexity, cooling requirements, and power consumption.

\subsubsection{Real-Time Processing and Data Transmission}
Processing 9 features per pixel is over $21\times$ faster than 192 features for KNN classification, with proportional memory savings.
For embedded systems, UAV-mounted sensors, or clinical point-of-care devices, this reduction enables practical deployment.
In remote sensing with bandwidth constraints, selective acquisition of only informative wavelengths could increase effective analysis rates.

\subsubsection{Training Data Efficiency}
Because wavelength selection is unsupervised (the autoencoder requires no labels), the framework could identify informative wavelengths from unlabeled data.
Subsequent supervised training for classification then operates on the reduced feature set, potentially requiring fewer labeled examples for equivalent performance.
This two-stage approach, unsupervised selection followed by supervised classification, is particularly valuable when labeled data is scarce or expensive to obtain.

% ----------------------------------------------------------------------------
\subsection{Comparison to Alternative Approaches}
\label{sec:discussion_alternatives}

Direct comparison to existing methods is complicated by the novelty of 4D ME-HSI methodology.
Most band selection literature addresses 3D HSI data from remote sensing; methods designed for EEM data focus on spectroscopic (non-imaging) applications.

For 3D HSI band selection, deep learning approaches include attention-based networks~\cite{mou2017deep} and reinforcement learning approaches~\cite{roy2020hybridsn}.
These methods require labeled training data and optimize selection for specific classification tasks.
The CAE perturbation-based unsupervised approach traded task-specific optimization for generalizability: the same wavelength selection could serve multiple downstream tasks (classification, segmentation, anomaly detection) without retraining.

For EEM spectroscopy, parallel factor analysis (PARAFAC)~\cite{bro1997parafac} decomposes 3D EEM tensors into component spectra, enabling chemical identification.
Extending PARAFAC to 4D imaging data with spatial dimensions is computationally challenging and loses spatial context that the convolutional approach preserved.

To the best of our knowledge, no prior work directly addressed wavelength selection for 4D ME-HSI data using deep learning approaches.


% ----------------------------------------------------------------------------
\subsection{Future Directions}
\label{sec:future_work}

The CAE perturbation-based framework provided a baseline for future research in this emerging area, with the potential for extension through attention mechanisms, variational formulations, or end-to-end optimization with task-specific losses. Specifically including: 
\subsubsection{Multi-Dataset Generalization}
Evaluation on a single lichen fluorescence dataset demonstrates
proof-of-concept but leaves cross-domain applicability as an open question.
Validating this framework on different datasets, each with distinct spectral
characteristics and class structures, would establish broader
generalizability. Transfer learning between related sample types could
further reduce the per-dataset training overhead.

\subsubsection{Automated Band Count Selection}
The current framework requires a user-specified target band count~$K$.
An adaptive selection criterion, such as monitoring marginal information
gain per additional band or by applying an elbow-point detector on the
accuracy, reduction curve, could
automate this choice, making the pipeline fully autonomous without
requiring domain knowledge or cross-validation.

\subsubsection{Interpretable Attribution}
Integrating domain-specific spectroscopic knowledge by including known fluorophore spectra with the perturbation-based importance scores could bridge the gap between statistical wavelength ranking and physical understanding. Such a hybrid attribution would increase practitioner trust and enable targeted experimental validation of the selected bands.

\subsubsection{Uncertainty-Aware Selection}
Replacing the deterministic autoencoder with a variational or Bayesian
formulations could yield confidence intervals on wavelength importance.
This is particularly valuable when selection guides expensive sensor design
decisions, where incorrect band choices carry a high cost.